\documentclass[11pt]{letter}
\usepackage[margin=2.5cm]{geometry}
\usepackage[T1]{fontenc}
\usepackage{hyperref}
\signature{Onishi Tatsuki\\Data Science AI Innovation Research Promotion Center, Shiga University\\bougtoir@gmail.com}
\address{Onishi Tatsuki\\Data Science AI Innovation Research Promotion Center, Shiga University\\1-1-1 Bamba, Hikone, Shiga 522-8522, Japan\\bougtoir@gmail.com}
\begin{document}
\begin{letter}{Prof.\ Claude Diebolt\\Managing Editor, \emph{Cliometrica}}
\opening{Dear Professor Diebolt,}

I am pleased to submit the manuscript entitled ``Network Exclusion and State Collapse: From Maritime Isolation to Technological Access Denial in the Long Run of History'' for consideration by \emph{Cliometrica}.

Using a comparative dataset of 96 historical polities spanning antiquity to the present, we distinguish between deliberate closure (policy-based trade bans, \emph{sakoku}, bloc membership) and what we term \emph{technical network exclusion}: involuntary disconnection from the dominant exchange network of an era due to geographic or technological constraints. The manuscript presents four features that we believe are of interest to the journal's readership.

First, a systematic sensitivity analysis demonstrates that reclassifying seven technically excluded polities transforms a non-significant association between closure and conquest (Fisher's exact $p = 0.187$) into a significant one ($p = 0.020$), while the core stock--flow odds ratio remains stable (OR $= 1.774$).

Second, we analyze conditional closure---cases where polities maintained selective technology channels while restricting broader trade. Examples include Tokugawa Japan's \emph{rangaku} through Dejima, Qing China's Canton system, Joseon Korea's tributary trade, and the Soviet Union's bloc-internal technology sharing. The mixed outcomes of these cases suggest that the open/closed dichotomy is insufficient; the conditions under which selective channels mitigate technology gaps merit further investigation.

Third, the dose--response gradient across closure types---technical exclusion (100\% conquest) $>$ policy bans $>$ \emph{sakoku} (with partial conduit) $>$ bloc $>$ open---points to technology flow disruption, rather than trade restriction per se, as the critical mechanism.

Fourth, the mechanism generalizes beyond geographic isolation: as the dominant network shifts from sea lanes to semiconductors, AI, and advanced robotics, states structurally excluded from these platforms may face analogous vulnerabilities.

The manuscript contains approximately 9{,}000 words, 4 tables, and 4 figures. Supplementary Table~S1 provides the complete dataset. We confirm that this work has not been published or submitted elsewhere.

We suggest the following reviewers based on their expertise in quantitative economic history:
\begin{itemize}
\item Prof.\ J\"org Baten (University of T\"ubingen) --- cliometric methods and long-run development
\item Prof.\ Stephen Broadberry (University of Oxford) --- comparative economic history and GDP estimation
\item Prof.\ Peter Turchin (Complexity Science Hub Vienna) --- quantitative historical dynamics and empire formation
\end{itemize}

\closing{Sincerely,}
\end{letter}
\end{document}